\chapter{Dysphagia (Difficulty Swallowing)}

\section{Definition}
Dysphagia is the subjective sensation of difficulty or abnormality in swallowing. It can involve the oral, pharyngeal, or esophageal phase of swallowing. Dysphagia affects approximately 15 percent of older adults and can lead to malnutrition, dehydration, and aspiration pneumonia.

\section{Pathophysiology}
Swallowing is a complex coordinated process involving cranial nerves V, VII, IX, X, and XII, as well as over 30 muscles. The oral phase involves mastication and bolus formation. The pharyngeal phase triggers the swallowing reflex, closing the airway. The esophageal phase uses peristalsis to move the bolus into the stomach. Dysphagia results from disruption at any phase.

\section{Etiology}
\begin{itemize}
\item Stroke (most common neurological cause)
\item Gastroesophageal reflux disease (GERD) with esophagitis
\item Esophageal stricture
\item Esophageal cancer
\item Neuromuscular disorders (Parkinson disease, multiple sclerosis, ALS)
\item Zenker diverticulum
\item Cricopharyngeal muscle dysfunction
\item Eosinophilic esophagitis
\item Achalasia (failure of LES relaxation)
\item Head and neck cancer or its treatment
\item Dementia or cognitive impairment
\end{itemize}

\section{Indications for Evaluation}
\begin{itemize}
\item Sensation of food sticking in the throat or chest
\item Pain with swallowing (odynophagia)
\item Coughing or choking during meals
\item Regurgitation of food or liquids
\item Recurrent pneumonia (from aspiration)
\item Unexplained weight loss or dehydration
\item Changes in voice quality after swallowing
\item Prolonged meal times or avoidance of certain foods
\end{itemize}

\section{Related Symptoms}
\begin{itemize}
\item Gastroesophageal reflux disease (GERD)
\item Heartburn
\item Regurgitation
\item Weight loss
\item Cough
\item Aspiration pneumonia
\item Sialorrhea (drooling)
\item Sore throat
\end{itemize}

\section{Self-Care/Home Management}
\begin{itemize}
\item Eat slowly and take small bites
\item Chew thoroughly before swallowing
\item Sit upright during and for 30 minutes after meals
\item Consider thickened liquids if thin liquids cause coughing
\item Moisten dry foods with sauce, gravy, or broth
\item Alternate bites of food with sips of liquid
\item Avoid talking while eating
\end{itemize}

\section{Diagnostic Approach}
\begin{itemize}
\item Bedside swallowing assessment by speech-language pathologist
\item Videofluoroscopic swallow study (modified barium swallow)
\item Fiberoptic endoscopic evaluation of swallowing (FEES)
\item Esophagogastroduodenoscopy (EGD) with biopsy
\item Esophageal manometry to assess motility
\item Barium swallow (esophagram)
\end{itemize}

\section{Treatment/Management Overview}
\begin{itemize}
\item Swallowing therapy by speech-language pathologist
\item Dietary modifications (texture-modified foods, thickened liquids)
\item Treat underlying cause (PPI for reflux, dilation for stricture)
\item Esophageal dilation for strictures
\item Botulinum toxin for cricopharyngeal dysfunction or achalasia
\item Surgical myotomy for achalasia
\item Percutaneous endoscopic gastrostomy (PEG) for severe dysphagia
\end{itemize}
